\section*{Capítulo \ref{ch:b2:plant-plasticity} --- Adaptaciones de las plantas y plasticidad fenotípica}

\begin{solution}{exo:b2:plant-plasticity:1}
Plasticidad: un \omterm{def:b2:meiosis-heredity:vocabulary}{genotipo}, varios \omterm{def:b2:meiosis-heredity:vocabulary}{fenotipos} según el ambiente (\omterm{prop:b2:plant-plasticity:sunshade}{hojas de sol}
y de sombra en un mismo roble). \omterm{def:b2:plant-plasticity:plasticity}{Norma de reacción}: la curva del \omterm{def:b2:meiosis-heredity:vocabulary}{fenotipo}
de un \omterm{def:b2:meiosis-heredity:vocabulary}{genotipo} en función de una variable ambiental (grosor de la hoja en
función de la luz). \omterm{def:b2:plant-plasticity:plasticity}{Aclimatación}: un ajuste fisiológico reversible (el
endurecimiento del centeno frente a las heladas en otoño). \omterm{def:b2:plant-plasticity:plasticity}{Adaptación}: un
ajuste heredable seleccionado a lo largo de generaciones (el metabolismo
CAM de un cactus).
\end{solution}

\begin{solution}{exo:b2:plant-plasticity:2}
Coleóptilo intacto iluminado por un lado: se curva hacia la luz. Extremo
cubierto con un capuchón opaco: no se curva --- el extremo percibe la luz.
Extremo cubierto con un capuchón transparente: se curva --- el capuchón en
sí es inocuo. Base protegida con un tubo opaco y extremo al descubierto:
se curva --- la percepción está en el extremo y la respuesta, por debajo.
\end{solution}

\begin{solution}{exo:b2:plant-plasticity:3}
\omterm{prop:b2:plant-plasticity:sunshade}{Hoja de sol}: pequeña, gruesa, dos o tres capas en empalizada, cutícula
gruesa, muchos estomas, fotosíntesis máxima elevada y punto de
compensación alto. \omterm{prop:b2:plant-plasticity:sunshade}{Hoja de sombra}: grande, fina, una capa en empalizada,
más clorofila por gramo, respiración baja, punto de compensación bajo. La
elección se hace cuando la hoja es un primordio en la yema, a partir de la
luz que recibe; una hoja madura no puede cambiar.
\end{solution}

\begin{solution}{exo:b2:plant-plasticity:4}
Las \omterm{prop:b2:plant-plasticity:stomata}{células oclusivas} captan potasio y agua, se hinchan y, como sus
paredes internas son más gruesas, se arquean separándose y abren el poro.
La luz azul y un $\mathrm{CO_2}$ interno bajo los abren por la mañana; el
\omterm{def:b2:plant-flowering:dormancy}{ácido abscísico}, procedente de las raíces en un suelo que se seca o de la
propia hoja, los cierra durante una sequía, igual que un aire muy seco.
\end{solution}

\begin{solution}{exo:b2:plant-plasticity:5}
$\varphi = R/(R + 0.77\,FR)$ con $FR = 1$: $1.15$: $0.60$; $0.7$:
$0.48$; $0.2$: $0.21$; $0.05$: $0.06$. El umbral de $0.4$ cae entre el
claro ($0.48$) y el dosel ($0.21$): la \omterm{thm:b2:plant-plasticity:phytochrome}{evitación de la sombra} se activa
bajo el dosel.
\end{solution}

\begin{solution}{exo:b2:plant-plasticity:6}
$E = 0.25\times 0.025 = \qty{6.25}{mmol.m^{-2}.s^{-1}}$; $A =
0.25\times 100/1.6 = \qty{15.6}{\micro mol.m^{-2}.s^{-1}}$; $A/E =
2.5\times 10^{-3}$ mol de carbono por mol de agua, es decir, 400 moléculas
de agua por carbono, $400\times 18/12 = \qty{600}{g}$ de agua por gramo de
carbono. En 12 horas: $6.25\times 10^{-3}\times 43200 =
\qty{270}{mol/m^2}$, \qty{4.9}{L} por metro cuadrado.
\end{solution}

\begin{solution}{exo:b2:plant-plasticity:7}
$E = 0.25\times 0.005 = \qty{1.25}{mmol.m^{-2}.s^{-1}}$, con $A$
sin cambios: $A/E = 0.0125$, 80 moléculas de agua por carbono,
\qty{120}{g/g} --- la quinta parte del coste diurno.
\end{solution}

\begin{solution}{exo:b2:plant-plasticity:8}
$v = 2r^{2}\Delta\rho g/9\eta = 2\times 4\times 10^{-12}\times
400\times 9.81/0.09 = \qty{3.5e-7}{m/s}$: \qty{0.35}{\micro m/s},
unos \qty{30}{s} para atravesar la célula. En un clinostato, la dirección
de la gravedad respecto a la célula cambia en cuestión de unos pocos
segundos de sedimentación, de modo que los granos nunca se acumulan en un
lado y el estímulo promediado a lo largo del tiempo de integración de la
planta (minutos) es nulo.
\end{solution}

\begin{solution}{exo:b2:plant-plasticity:9}
Alargamiento medio en dos horas, $0.04$; lados, $1.2\times 0.04 = 0.048$
y $0.8\times 0.04 = 0.032$; $\theta = 20\times 0.016/2 = 0.16$
rad, unos \qty{9}{\degree}. La curvatura se acumula a razón de $0.08$ rad
por hora; $\pi/2$ requiere unas \qty{20}{h}.
\end{solution}

\begin{solution}{exo:b2:plant-plasticity:10}
Las reservas permiten \qty{50}{cm} de tallo. Ortigal: \qty{30}{cm} en
7,5 días por \qty{30}{mg} --- la plántula alcanza la luz con reservas de
sobra. Bosque: \qty{10}{m} están fuera de su alcance; la plántula gasta sus
\qty{50}{mg} en medio metro de tallo pálido en doce días y muere
ahilada, mientras que una estrategia de tolerancia a la sombra --- unas
pocas hojas finas, un crecimiento lento --- podría haberla mantenido viva
durante años.
\end{solution}

\begin{solution}{exo:b2:plant-plasticity:11}
El hielo se forma primero en los espacios extracelulares, donde la
solución está más diluida; la presión de vapor sobre el hielo es menor que
sobre el agua de la célula, de modo que el agua sale de la célula para
unirse al hielo, y la célula se deshidrata y sus membranas se colapsan.
Una semana de frío carga la célula de azúcares y proteínas protectoras que
bajan su punto de congelación y estabilizan las membranas, y fabrica
proteínas anticongelantes que limitan el crecimiento de los cristales.
Mantener el programa activo todo el año desvía carbono del crecimiento a
la protección y frena la división: una planta permanentemente endurecida
es una planta pequeña.
\end{solution}

\begin{solution}{exo:b2:plant-plasticity:12}
La señal del rojo lejano es honesta cuando la vecina puede ser superada y
mentirosa cuando es un árbol; la semana cálida es honesta en abril y
mentirosa en febrero; en un cultivo denso, las vecinas de cada planta son
sus iguales, de modo que la señal es honesta en su forma, pero inútil, y
los mejoradores seleccionan plantas que la ignoran. Una \omterm{def:b2:plant-plasticity:plasticity}{norma de reacción}
es el registro de la frecuencia con que cada señal ha dicho la verdad en
el pasado del linaje, y falla siempre que el presente difiere de ese
pasado.
\end{solution}

\begin{solution}{pb:b2:plant-plasticity:1}
\textbf{1.} Al descubierto: $1.15/(1.15 + 0.77) = 0.60$; bajo el dosel:
$0.2/(0.2 +
0.77) = 0.21$.
\textbf{2.} $\varphi$ bajo: plantas encima --- alargarse. Bajo una malla de
sombreo, $\varphi$ se mantiene en $0.60$: un día oscuro, no una vecina; no
hay \omterm{thm:b2:plant-plasticity:phytochrome}{evitación de la sombra}, solo un crecimiento más lento por falta de
luz.
\textbf{3.} $1 + 2(0.60 - 0.21) = 1.78$: \qty{2.7}{cm} al día.
\textbf{4.} \qty{37}{cm} frente a \qty{21}{cm}.
\textbf{5.} La curvatura hacia el claro (fototropismo, a través del
receptor de luz azul, la fototropina) y un alargamiento más rápido y una
reorientación de las hojas en el lado con mayor $R{:}FR$ (a través del
fitocromo).
\textbf{6.} Las \omterm{prop:b2:plant-plasticity:sunshade}{hojas de sombra} se construyen grandes y finas para captar
de forma barata una luz escasa; a pleno sol se sobrecalientan, sufren
fotoinhibición y se queman, y la planta debe sustituirlas por \omterm{prop:b2:plant-plasticity:sunshade}{hojas de
sol}.
\textbf{7.} \qty{2.7}{cm} al día a lo largo de una zona de \qty{15}{mm} son
\qty{0.11}{cm} por hora, un alargamiento relativo de $0.074$ por hora;
lado sombreado, $1.3\times 0.074 = 0.096$; lado iluminado, $0.7\times 0.074 =
0.052$.
\textbf{8.} $\theta = 15\times(0.096 - 0.052)/2 = 0.33$ rad, unos
\qty{19}{\degree} en una hora.
\textbf{9.} \qty{60}{\degree} son $1.05$ rad: unas tres horas.
\textbf{10.} $v = 2\times 6.25\times 10^{-12}\times 400\times
9.81/0.09 = \qty{5.5e-7}{m/s}$, \qty{0.55}{\micro m/s}: \qty{22}{s}
para atravesar la célula.
\textbf{11.} El fototropismo, y la propia \omterm{thm:b2:plant-plasticity:phytochrome}{evitación de la sombra} eleva el
ángulo con el que el tallo se conforma con crecer: la luz es el recurso
por el que se compite, y un tallo que se enderezara contra la gravedad
volvería hacia la sombra.
\textbf{12.} La raíz se curva hacia abajo hasta quedar vertical: los
estatolitos sedimentan en el lado inferior, la auxina se acumula allí y,
en una raíz, el exceso de auxina inhibe el crecimiento, de modo que el
lado superior crece más deprisa. La misma redistribución, con una
sensibilidad opuesta de las células.
\textbf{13.} $E = 0.2\times 0.015 = \qty{3}{mmol.m^{-2}.s^{-1}}$; sobre
$2\times 10^{-3}$ \unit{m^2} y durante \qty{50400}{s}: \qty{0.30}{mol},
\qty{5.4}{g} de agua al día.
\textbf{14.} $A = 0.2\times 80/1.6 = \qty{10}{\micro mol.m^{-2}.s^{-1}}$;
al día, $0.50$ mol de carbono por metro cuadrado; $A/E = 3.3\times
10^{-3}$: 300 moléculas de agua por carbono, \qty{450}{g/g}.
\textbf{15.} $0.50\times 2\times 10^{-3} = \qty{1}{mmol}$: \qty{12}{mg}
de carbono al día.
\textbf{16.} Claro: $E = \qty{6}{mmol.m^{-2}.s^{-1}}$, \qty{10.9}{g} al
día; eficiencia de \qty{900}{g/g}.
\textbf{17.} $E = 0.02\times 0.03 = \qty{0.6}{mmol.m^{-2}.s^{-1}}$
(diez veces menos); $A = 0.02\times 200/1.6 = \qty{2.5}{\micro
mol.m^{-2}.s^{-1}}$ (cuatro veces menos). El agua disminuye más: al
reducir la hoja su $\mathrm{CO_2}$ interno, el gradiente para el carbono
se amplía, mientras que el del agua no puede hacerlo.
\textbf{18.} \qty{4.25}{g} de agua, de los que puede perder \qty{0.85}{g};
a \qty{10.9}{g} al día ($0.78$ g por hora) se marchita en una hora
aproximadamente.
\textbf{19.} El CAM es lento y necesita tejido suculento de reserva; una
plántula que compite por la luz bajo un dosel, donde la pérdida de agua es
de todos modos pequeña, necesita rapidez, no economía.
\textbf{20.} $100/2 = \qty{50}{cm}$.
\textbf{21.} Ortigal: \qty{40}{cm} en 15 días por \qty{80}{mg} --- alcanza
la luz. Bosque: \qty{8}{m} son inalcanzables; gasta sus reservas en medio
metro y muere. Mejor, en el bosque: quedarse baja, formar hojas finas de
sombra y esperar a que se abra un claro --- la manera del haya.
\textbf{22.} Abedul: una norma empinada, con un fuerte alargamiento a
$\varphi$ bajo, puesto que las vecinas de un pionero son de su tamaño.
Haya: una norma plana, con poco alargamiento y mucha tolerancia, puesto
que sus vecinos son árboles.
\textbf{23.} Desactivar la respuesta de \omterm{thm:b2:plant-plasticity:phytochrome}{evitación de la sombra} ---
mantener activo el fitocromo B o eliminar los \omterm{prop:b2:cell-differentiation:transcriptional}{factores de transcripción}
que este frena --- y comprobar la época de floración, la germinación y la
resistencia del tallo, que controla la misma vía.
\textbf{24.} Una planta que no puede percibir a sus vecinas queda superada
siempre que habría podido ganar; una planta que siempre cree la señal se
malgasta bajo cada árbol. La norma ganadora apuesta por la señal en
proporción a la frecuencia con que, en el pasado del linaje, acertó.
\textbf{25.} $\varphi = 0.60$ al descubierto y $0.21$ bajo el dosel;
\qty{37}{cm} al cabo de 14 días; \qty{5.4}{g} de agua al día a la sombra y
\qty{10.9}{g} en el claro.
\end{solution}
